\section{Cognitive Map: el espacio es solo un caso particular del razonamiento estructural}

El mapa cognitivo surgió originalmente de la navigation, pero el objetivo del curso es extenderlo a cualquier relational structure. Primero revisamos la evidencia espacial, después la transitive inference no espacial y, por último, abstraemos el problema común como graph learning: el contenido observado cambia, las relaciones entre aristas se mantienen y el modelo necesita reutilizar la estructura y actualizar con rapidez la vinculación del contenido.

\subsection{Del mapa de Tolman a la generalización en árbol}

La diapositiva Theory/navigation describe el cognitive map como un modelo interno de la estructura del entorno, que permite al agent replanificar cuando un camino se bloquea o cambia la meta, en lugar de limitarse a reproducir una cached action sequence. El ejemplo en árbol muestra además que la estructura no tiene que ser un plano bidimensional: las jerarquías, las relaciones de parentesco y los estados de una tarea también pueden formar un graph; lo decisivo es si el modelo puede inferir trayectorias que no ha experimentado directamente.

\lecturefigure{slide-43-cognitive-map-theory-navigation.jpg}{Teoría del cognitive map: un modelo estructural permite una navegación flexible}{Estado de enseñanza manuscrito del video oficial V043; 00:47:32}

\lecturefigure{slide-44-tree-structured-generalization.jpg}{Generalización estructural en relaciones en forma de árbol}{Estado de enseñanza manuscrito del video oficial V044; 00:48:20}

Si el entorno se representa como un grafo $G=(V,E)$, cada nodo $i$ tiene una observación $x_i$ y el tipo de arista o acción $a_t$ determina la transition $i_t\to i_{t+1}$. La generalización estructural exige aprender
$$
g_{t+1}=f_{\theta}(g_t,a_t),
$$
donde $g_t$ es un structural state lo más desacoplado posible de la observation identity concreta. Al cambiar el conjunto de $x_i$, $f_{\theta}$ sigue pudiendo actualizar la posición o las relaciones; solo hace falta volver a vincular $g_i$ con el nuevo contenido.

\begin{knowledgebox}{Concepto previo: un cognitive map no equivale a un mapa geográfico}
Lo esencial de un mapa es que las relaciones se pueden componer: si A está a la izquierda de B y B a la izquierda de C, se deduce la posición relativa de A y C; si cierta acción lleva del nodo 1 al nodo 2, la misma regla de acción puede transferirse a objetos nuevos. El espacio bidimensional es el caso más intuitivo, pero los graph, los family tree, los estados de una tarea y las relaciones entre conceptos pueden usar la misma idea estructural.
\end{knowledgebox}

\subsection{Place cells, grid cells y evidencia de estructura espacial}

La diapositiva Evidence overview reúne tareas espaciales, neural recording y resultados conductuales. El ponente repasa de palabra la place cell: cierta hippocampal neuron presenta un firing intenso cuando el animal se encuentra en una ubicación específica del entorno; la grid cell, en cambio, forma una hexagonal lattice sobre múltiples ubicaciones. No son coordenadas GPS directas, sino que la actividad neuronal muestra una estructura estable en función de la posición.

\lecturefigure{slide-45-task-evidence-selection.jpg}{Panorama de la evidencia conductual y neuronal del mapa cognitivo}{Estado de enseñanza manuscrito del video oficial V045; 00:51:44}

\teachervoice{Will usa el examen «Knowledge» de los London taxi drivers para ilustrar el costo conductual de aprender una estructura espacial, y cita la evidencia de correlación entre el hippocampal size y la experiencia de manejo. Enseguida pasa a tareas no espaciales, porque una «correlación espacial» todavía no demuestra que la región cerebral represente una relational structure más general.}

\begin{warningbox}{Interpretaciones erróneas frecuentes de las tuning curve neuronales}
Que una cell presente firing en una ubicación específica no significa que ella sola «almacene esa ubicación»; la posición suele representarse mediante un population code. Que el grid pattern se parezca a la position encoding de un modelo tampoco implica que el proceso de entrenamiento, el origen de las entradas y la función biológica sean idénticos. Lo que debe compararse es la información decodificable, las regularidades de transformación y el papel causal en la tarea.
\end{warningbox}

\subsection{Evidencia no espacial: Transitive inference}

La diapositiva Transitive-inference presenta un tipo de tarea relacional no espacial: el animal aprende pair locales como A mayor que B, B mayor que C, etc., y después se le evalúa con A y C, que no se entrenaron directamente. Si solo memoriza los observed pairs, no puede inferir de forma estable las combinaciones nuevas; si forma un ordered relational map, puede deducirlas siguiendo la estructura. La diapositiva manuscrita coloca los experimentos conductuales, la evidencia de imagen y la pair relation en una misma cadena argumentativa.

\lecturefigure{slide-46-transitive-inference-evidence.jpg}{Transitive inference: inferir relaciones no vistas a partir de pair aprendidos}{Estado de enseñanza manuscrito del video oficial V046; 00:54:54}

\lecturefigure{slide-47-executive-layer-map.jpg}{Observation, structure y lectura conductual en tareas relacionales}{Estado de enseñanza manuscrito del video oficial V047; 00:55:18}

\begin{knowledgebox}{Lectura de la figura: cómo distinguir las tres estrategias}
\term{pair memorization} solo puede responder sobre los pair entrenados; \term{scalar ordering} asigna a cada objeto un rank unidimensional, adecuado para órdenes lineales; \term{relational graph} representa varios tipos de edge y puede manejar árboles, ciclos y espacios bidimensionales. El diseño experimental debe distinguir estas estrategias mediante combinaciones no vistas y cambios de estructura, en lugar de medir solo la pair accuracy en entrenamiento.
\end{knowledgebox}

\subsection{Del graph world model a un algoritmo neuronal general}

La diapositiva Graph-world abstrae el entorno como nodos, aristas y observation; la siguiente enumera el algoritmo objetivo: la estructura debe poder reutilizarse entre entornos, el contenido debe escribirse con rapidez, la planificación debe ocurrir sobre el latent graph y las representaciones neuronales deben contrastarse con datos hippocampal/entorhinal. El reto abarca a la vez representation learning, memory binding y model-based inference.

\lecturefigure{slide-48-graph-world-model.jpg}{Abstraer el entorno como un graph reutilizable y observations variables}{Estado de enseñanza manuscrito del video oficial V048; 00:58:58}

\lecturefigure{slide-49-neural-algorithm-general.jpg}{Algoritmo neuronal ideal: transferencia estructural, vinculación rápida e inferencia}{Estado de enseñanza manuscrito del video oficial V049; 00:59:28}

\begin{importantbox}{Descomposición de la tarea en la segunda clase}
El aprendizaje lento se encarga de dominar la transition structure y las reglas componibles; el aprendizaje rápido se encarga de vincular el nuevo sensory content con los structural states; la inferencia se encarga de propagar y planificar a lo largo de la estructura. La arquitectura de TEM asigna precisamente estas tres funciones a representaciones y mecanismos de actualización distintos.
\end{importantbox}

\subsection{Resumen de la sección}

Las place/grid cells y la taxi-driver evidence respaldan el mapa cognitivo espacial, mientras que tareas como la transitive inference respaldan un relational map más amplio. El problema computacional común es separar la observation identity de la reusable structure y después vincularlas mediante memoria rápida, de modo que también puedan inferirse combinaciones no vistas.
